% Gerado de docs/monografia/apendice-protocolo.md por scripts/monografia_latex.py.
% Não edite aqui: edite o .md e gere de novo.

\chapter{\texorpdfstring{\mbox{PROTOCOLO} \mbox{DO} \mbox{ESTUDO} \mbox{COM} \mbox{ESTUDANTES}}{PROTOCOLO DO ESTUDO COM ESTUDANTES}}\label{cap:A}

O estudo que este trabalho deixa para depois (7.2), desenhado para a versão do sistema descrita no Capítulo 3. Adaptado de \texttt{docs/\allowbreak{}plano-{}tcc.\allowbreak{}md}, que foi escrito para a versão anterior da sessão.

\section{As quatro perguntas}

A pergunta do estudo é se é viável diagnosticar concepções errôneas sobre ponteiros em C a partir da resposta, da razão e da certeza declarada, e confrontá-las com demonstrações, numa sessão que estudantes de graduação consigam e queiram completar. Ela se divide em quatro:

\begin{quadro}[htbp]
\caption{As quatro perguntas do estudo}\label{qua:A-1}
\footnotesize
\begin{tabularx}{\textwidth}{|>{\centering\arraybackslash}p{1.2cm}|X|X|X|}
\hline
\textbf{} & \textbf{Pergunta} & \textbf{Medida} & \textbf{De onde vem} \\ \hline
V1 & Os estudantes concluem a sessão sem ajuda de conteúdo do pesquisador, e em quanto tempo? & conclui (sim/não); minutos por tópico; onde parou & o registro da sessão \\ \hline
V2 & A concepção errônea detectada é compatível com a explicação que o participante dá pensando em voz alta? & compatível / incompatível / inconclusiva, por episódio & a gravação de áudio × as crenças registradas na sessão \\ \hline
V3 & Depois da demonstração, o participante responde de outro jeito? & a reconferência do fim: certa pela razão certa, ou ``a ideia voltou'' & o registro da sessão \\ \hline
V4 & Como percebe clareza, dificuldade e utilidade? & SUS, de 10 itens \cite{brooke1996}, e três perguntas abertas & questionário ao fim \\ \hline
\end{tabularx}
\fonte{Elaborado pelo autor.}
\end{quadro}

A unidade de análise de V2 e V3 \textbf{não é o participante}: é o episódio de concepção errônea detectada. Dez pessoas podem gerar 25 episódios ou 4, e o denominador de cada resultado é esse número.

\section{A sessão}

De 10 a 15 estudantes que cursam ou já cursaram a disciplina com C, um de cada vez, presencialmente, em cerca de 30 minutos cada:

\begin{enumerate}
\item Consentimento e código anônimo (P01…P15). Nome e matrícula ficam em papel, fora do sistema.
\item Três perguntas rápidas sobre experiência com C.
\item Dois tópicos na trilha, pensando em voz alta. O pesquisador usa apenas lembretes padronizados (``continue pensando em voz alta''), definidos antes da coleta: nenhuma pista sobre o conteúdo, nenhuma correção. Sugestão de tópicos: kc\_07 (ponteiro e array) e kc\_09 (\texttt{free} e ponteiro pendente), os que têm mais concepções mapeadas.
\item Entrevista curta: o que achou, o que confundiu, faria outro tópico.
\end{enumerate}

\section{A classificação dos episódios de V2}

\begin{quadro}[htbp]
\caption{A classificação dos episódios de V2}\label{qua:A-2}
\footnotesize
\begin{tabularx}{\textwidth}{|X|X|}
\hline
\textbf{Classificação} & \textbf{Critério} \\ \hline
Compatível & a fala traz explicitamente a mesma relação conceitual que a concepção afirma \\ \hline
Incompatível & a fala apresenta explicação que contradiz a concepção detectada (inclusive ``li errado o enunciado'') \\ \hline
Inconclusiva & a fala não traz evidência suficiente para decidir \\ \hline
\end{tabularx}
\fonte{Elaborado pelo autor.}
\end{quadro}

Uma segunda pessoa classifica pelo menos um terço dos episódios, às cegas. Com poucos episódios, reporta-se a contagem e os casos, e não um coeficiente de concordância. Com 10 a 15 pessoas não há teste de hipótese: os números são descritivos, e o que sustenta as conclusões são os casos: um episódio incompatível vale mais do que dez compatíveis, porque mostra onde a regra falha.

\section{Ameaças à validade}

\begin{quadro}[htbp]
\caption{Ameaças à validade e contramedidas}\label{qua:A-3}
\footnotesize
\begin{tabularx}{\textwidth}{|X|X|}
\hline
\textbf{Ameaça} & \textbf{Contramedida} \\ \hline
A reconferência vem logo depois da demonstração & ela vem no fim do tópico, com as outras premissas no meio \\ \hline
Ninguém erra com convicção, e V2 e V3 ficam sem episódios & recrutar quem está cursando a disciplina; tópicos com mais concepções mapeadas; medir a taxa num ensaio com dois colegas antes de marcar a coleta \\ \hline
Pensar em voz alta muda o próprio raciocínio & tratado como evidência qualitativa, nunca como leitura do que a pessoa pensa \\ \hline
O pesquisador contamina a sessão & lembretes padronizados, escritos antes \\ \hline
A hipótese está certa, mas a fala não revela & a categoria ``inconclusiva'' existe e é reportada \\ \hline
\end{tabularx}
\fonte{Elaborado pelo autor.}
\end{quadro}

\section{O que falta antes da coleta}

\begin{enumerate}
\item A exportação das sessões da trilha em tabelas (hoje a exportação lê só a versão anterior da sessão).
\item O termo de consentimento, e a submissão ao comitê de ética, se exigida.
\item O ensaio com dois colegas.
\item Um procedimento de exclusão a pedido que alcance também as cópias de segurança do registro.
\end{enumerate}
